\documentclass[a4paper,10pt]{report}\input{../0.Préambule/Préambule_cours_prof}\begin{document}

\newpage
\section*{Aires}
\addcontentsline{toc}{section}{Aires}

\subsection*{Aire et unités}
\addcontentsline{toc}{section}{Aire et unités}

\begin{exer1}
Complète.

\begin{enumerate}[font=\color{exer1}]
\renewcommand{\arraystretch}{1.2}
\begin{tabularx}{\linewidth}{m{1.7cm}*{1}{>{\arraybackslash}X}m{0.5cm}m{2cm}*{1}{>{\arraybackslash}X}}
\item $3$m$^2$ $= $ & \dotfill cm$^2$&&
\item $700$cm$^2$ $= $ & \dotfill m$^2$\\
\item $5$m$^2$ $= $ & \dotfill dm$^2$&&
\item $6000$dm$^2$ $= $ & \dotfill m$^2$\\
\item $5$cm$^2$ $= $ & \dotfill mm$^2$&&
\item $400$mm$^2$ $= $ & \dotfill cm$^2$\\
\end{tabularx}
\end{enumerate}
\end{exer1}

\begin{exer2}
On donne les superficies suivantes :
\begin{enumerate}[font=\color{exer2}]
\renewcommand{\arraystretch}{0.7}
\begin{tabularx}{\linewidth}{*{2}{>{\arraybackslash}X}}
\item[•] Belle-île-en-mer : $90$km$^2$&
\item[•] Belle-île-en-mer : $\np{2300}$ ha \\
\item[•] Belle-île-en-mer : $\np{175000000}$m$^2$&
\item[•] Belle-île-en-mer : $\np{1160000}$dam$^2$\\
\end{tabularx}
\end{enumerate}
Range ces iles dans l'ordre décroissant de leur superficie.
\end{exer2}

\subsection*{Calculer une aire}
\addcontentsline{toc}{section}{Calculer une aire}

\vspace{-1cm}
\noindent \begin{minipage}{10cm}
\begin{exer1}
Donne la nature de chaque quadrilatère en justifiant, puis calcule (sans calculatrice) son aire (pose si nécessaire les opérations)
\end{exer1}
\end{minipage}
\hspace{3mm}
\begin{minipage}{6cm}
\begin{center}
\includegraphics[scale=0.6]{../40.Image/14. Exo-Les aires 7.png}
\end{center}
\end{minipage}

\begin{exer1}
Calcule les aires de chacun des triangles. (Attention, les triangles ne sont pas dessinés en vraie grandeur.)

\begin{center}
\includegraphics[scale=0.65]{../40.Image/14. Exo-Les aires 10.png}
\end{center}
\end{exer1}

\begin{exer2}
$HSLE$ est un rectangle d'aire $18$cm$^2$ et $L=6cm$. Détermine la longueur du segment $[LS]$.
\end{exer2}

\begin{exer2}
$RSV$ est un triangle rectangle en $S$ d'aire $7$cm$^2$ et $RS=2cm$. Détermine la longueur de $[VS]$.
\end{exer2}

\begin{exer2}
$GLSM$ est un carrée. $[GS]$ mesure $3$cm. $GSB$ est un triangle rectangle en $S$, $B$ étant à l'extérieur du carré. $[BS]$ mesure $5$cm. Trace une figure a main levée, codée, puis calcule sont aire totale.
\end{exer2}

\begin{exer1}
Calcule l'aire d'un disque de rayon $4$cm et l'aire d'un disque de diamètre $12$cm. On donnera le calcul (avec $\pi$) et le résultat arrondie au dixième près.
\end{exer1}

\begin{exer1}
Calcule l'aire d'un disque de rayon $4$m et l'aire d'un disque de diamètre $4,3$hm. On donnera le calcul (avec $\pi$) et le résultat arrondie au centième près.
\end{exer1}

\noindent \begin{minipage}{9cm}
\begin{exer1}
Calcule l'ire des figures colorées suivantes. Donne le calcul (avec le nombre $\pi$), puis la valeur arrondie au centième.
\end{exer1}
\end{minipage}
\begin{minipage}{8cm}
\begin{center}
\includegraphics[scale=0.65]{../40.Image/14. Exo-Les aires 12.png}
\end{center}
\end{minipage}

\newpage
\subsection*{Aire et unités}
\addcontentsline{toc}{section}{Aire et unités}

\vspace{-1cm}
\noindent \begin{minipage}{11cm}
\begin{exer1}
Calcule l'aire du parallélogramme $ABCD$ ci-contre.
\end{exer1}
\end{minipage}
\begin{minipage}{6cm}
\begin{center}
\includegraphics[scale=0.5]{../40.Image/14. Exo-Les aires 13.png}
\end{center}
\end{minipage}

\begin{exer1}
Dans chaque cas, calcule l'aire du parallélogramme $ABCD$ en mesurant les longueurs nécessaires.

\begin{center}
\includegraphics[scale=0.4]{../40.Image/14. Exo-Les aires 14.png}
\end{center}
\end{exer1}

\vspace{-0.5cm}
\noindent \begin{minipage}{11.5cm}
\begin{exer1}
Calcule la longueur marquée par un point d'interrogation.
\end{exer1}
\end{minipage}
\begin{minipage}{5.5cm}
\begin{center}
\includegraphics[scale=0.3]{../40.Image/14. Exo-Les aires 15.png}
\end{center}
\end{minipage}

\noindent \begin{minipage}{10cm}
\begin{exer1}
On donne la figure ci-contre, qui représente un cerf-volant.

\medskip
\noindent A l'aide du codage et des informations fournies, calcule son aire.
\end{exer1}
\end{minipage}
\begin{minipage}{7cm}
\begin{center}
\includegraphics[scale=0.3]{../40.Image/14. Exo-Les aires 16.png}
\end{center}
\end{minipage}

\end{document}